\documentclass[12pt]{article}
\usepackage[UTF8]{inputenc}
\usepackage{thaisupp}
\usepackage{enumitem}
\begin{document}
\title{ความร้อน}
\maketitle
\section*{In-class Questions}
\begin{enumerate}[label=\arabic*.]
\item ความร้อน (Heat) กับ อุณหภูมิ (Temperature) แตกต่างกันอย่างไร \hfill \textbf{\small (Main)}
\begin{enumerate}[label=)]
\item ก) เป็นสิ่งเดียวกัน
\item ข) ความร้อนคือพลังงานที่ถ่ายโอน อุณหภูมิคือมาตรวัดระดับร้อน
\item ค) อุณหภูมิคือพลังงาน ความร้อนคือการวัด
\item ง) ความร้อนวัดเป็นองศา อุณหภูมิวัดเป็นจูล
\end{enumerate}
\item น้ำ 1 kg รับความร้อน 4,200 J อุณหภูมิจะเพิ่มขึ้นกี่เซลเซียส (ความจุความร้อนจำเพาะของน้ำ = 4,200 J/kg·°C) \hfill \textbf{\small (Main)}
\begin{enumerate}[label=)]
\item ก) 1 °C
\item ข) 4.2 °C
\item ค) 420 °C
\item ง) 0.42 °C
\end{enumerate}
\item การถ่ายโอนความร้อนมีกี่วิธีหลัก \hfill \textbf{\small (Main)}
\begin{enumerate}[label=)]
\item ก) 1 วิธี
\item ข) 2 วิธี
\item ค) 3 วิธี
\item ง) 4 วิธี
\end{enumerate}
\item สมการแก๊สอุดมคติ PV = nRT โดย R มีค่าเท่าไร \hfill \textbf{\small (Main)}
\begin{enumerate}[label=)]
\item ก) 8.314 J/(mol·K)
\item ข) 0.0821 L·atm/(mol·K)
\item ค) ทั้ง ก และ ข ถูกต้อง
\item ง) ขึ้นอยู่กับชนิดของแก๊ส
\end{enumerate}
\item กฎของบอยล์ (Boyle's Law) ระบุว่าอะไร เมื่ออุณหภูมิคงที่ \hfill \textbf{\small (Main)}
\begin{enumerate}[label=)]
\item ก) P ∝ V
\item ข) P ∝ 1/V
\item ค) P ∝ T
\item ง) V ∝ T
\end{enumerate}
\item กฎของชาร์ลส์ (Charles's Law) ระบุว่าอะไร เมื่อความดันคงที่ \hfill \textbf{\small (Main)}
\begin{enumerate}[label=)]
\item ก) P ∝ V
\item ข) P ∝ 1/V
\item ค) V ∝ T
\item ง) P ∝ T
\end{enumerate}
\item การนำความร้อน (Conduction) เกิดขึ้นได้ดีในสารใด \hfill \textbf{\small (Main)}
\begin{enumerate}[label=)]
\item ก) ของแข็ง โดยเฉพาะโลหะ
\item ข) ของเหลวเท่านั้น
\item ค) แก๊สเท่านั้น
\item ง) สุญญากาศเท่านั้น
\end{enumerate}
\item แก๊สอุดมคติ 1 mol ที่ STP (0°C, 1 atm) มีปริมาตรเท่าใด \hfill \textbf{\small (Main)}
\begin{enumerate}[label=)]
\item ก) 22.4 L
\item ข) 24.0 L
\item ค) 20.0 L
\item ง) 2.24 L
\end{enumerate}
\item เหล็กมวล 2 kg รับความร้อน 9,000 J อุณหภูมิเพิ่มขึ้นเท่าไร (c ของเหล็ก = 450 J/kg·°C) \hfill \textbf{\small (Main)}
\begin{enumerate}[label=)]
\item ก) 5 °C
\item ข) 10 °C
\item ค) 20 °C
\item ง) 1 °C
\end{enumerate}
\item การพาความร้อน (Convection) เกิดขึ้นได้ดีในสารใด \hfill \textbf{\small (Main)}
\begin{enumerate}[label=)]
\item ก) ของแข็ง
\item ข) ของเหลวและแก๊ส
\item ค) สุญญากาศ
\item ง) โลหะเท่านั้น
\end{enumerate}
\item ความดันของแก๊ส 2 atm มีค่าเท่ากับกี่ Pa \hfill \textbf{\small (Main)}
\begin{enumerate}[label=)]
\item ก) 2 × 10⁵ Pa
\item ข) 2 × 10⁴ Pa
\item ค) 1.013 × 10⁵ Pa
\item ง) 2.026 × 10⁵ Pa
\end{enumerate}
\item แก๊สอุดมคติ 2 mol ที่อุณหภูมิ 300 K และความดัน 1 × 10⁵ Pa มีปริมาตรเท่าใด (R = 8.314 J/mol·K) \hfill \textbf{\small (Main)}
\begin{enumerate}[label=)]
\item ก) 2.5 L
\item ข) 5.0 L
\item ค) 25 L
\item ง) 50 L
\end{enumerate}
\item ดวงอาทิตย์ส่งความร้อนมาถึงโลกโดยวิธีใด \hfill \textbf{\small (Main)}
\begin{enumerate}[label=)]
\item ก) การนำความร้อน
\item ข) การพาความร้อน
\item ค) การแผ่รังสี
\item ง) การพาความร้อนและแผ่รังสี
\end{enumerate}
\item น้ำแข็ง 0.5 kg ที่ 0°C ละลายเป็นน้ำที่ 0°C ต้องใช้ความร้อนเท่าไร (L ของน้ำแข็ง = 334,000 J/kg) \hfill \textbf{\small (Main)}
\begin{enumerate}[label=)]
\item ก) 167,000 J
\item ข) 334,000 J
\item ค) 668,000 J
\item ง) 83,500 J
\end{enumerate}
\item เมื่ออุณหภูมิของแก๊สเพิ่มขึ้น 2 เท่า (ปริมาตรคงที่) ความดันจะเป็นอย่างไร \hfill \textbf{\small (Main)}
\begin{enumerate}[label=)]
\item ก) ความดันลดลง 2 เท่า
\item ข) ความดันเพิ่มขึ้น 2 เท่า
\item ค) ความดันคงที่
\item ง) ความดันเพิ่มขึ้น 4 เท่า
\end{enumerate}
\item ข้อใดเกี่ยวกับการแผ่รังสีความร้อนได้ถูกต้อง \hfill \textbf{\small (Main)}
\begin{enumerate}[label=)]
\item ก) ต้องอาศัยตัวกลาง
\item ข) เกิดขึ้นเฉพาะในของแข็ง
\item ค) เป็นคลื่นแม่เหล็กไฟฟ้า
\item ง) เกิดขึ้นเฉพาะในแก๊ส
\end{enumerate}
\item ทองแดงมีความจุความร้อนจำเพาะ 390 J/kg·°C หมายความว่าอะไร \hfill \textbf{\small (Main)}
\begin{enumerate}[label=)]
\item ก) ต้องใช้ 390 J เพื่อทำให้ทองแดง 1 kg ร้อนขึ้น 1°C
\item ข) ทองแดง 1 kg สามารถเก็บความร้อนได้ 390 J
\item ค) ทองแดงนำความร้อนได้ดีกว่าน้ำ 10 เท่า
\item ง) ถูกทั้ง ก และ ข
\end{enumerate}
\item ถังแก๊สปริมาตร 10 L บรรจุแก๊สที่ความดัน 3 atm และอุณหภูมิ 300 K มีกี่โมล (R = 0.082 L·atm/mol·K) \hfill \textbf{\small (Main)}
\begin{enumerate}[label=)]
\item ก) 0.61 mol
\item ข) 1.22 mol
\item ค) 2.44 mol
\item ง) 0.12 mol
\end{enumerate}
\item น้ำ 500 g ที่ 20°C ได้รับความร้อนจนมีอุณหภูมิ 70°C ต้องใช้ความร้อนเท่าไร (c ของน้ำ = 4,200 J/kg·°C) \hfill \textbf{\small (Main)}
\begin{enumerate}[label=)]
\item ก) 105,000 J
\item ข) 210,000 J
\item ค) 420,000 J
\item ง) 525,000 J
\end{enumerate}
\item STP ย่อมาจากอะไร \hfill \textbf{\small (Main)}
\begin{enumerate}[label=)]
\item ก) Standard Temperature and Pressure
\item ข) Standard Time and Pressure
\item ค) Stable Temperature and Pressure
\item ง) Standard Thermal Process
\end{enumerate}
\item ช้อนโลหะจุ่มในน้ำร้อน ความร้อนถ่ายโอนผ่านช้อนโดยวิธีใด \hfill \textbf{\small (Main)}
\begin{enumerate}[label=)]
\item ก) การพาความร้อน
\item ข) การแผ่รังสี
\item ค) การนำความร้อน
\item ง) ทั้งสามวิธีพร้อมกัน
\end{enumerate}
\item เหล็ก 1 kg ที่ 100°C วางในน้ำ 2 kg ที่ 20°C อุณหภูมิสุดท้ายเป็นเท่าไร (c เหล็ก = 450 J/kg·°C, c น้ำ = 4,200 J/kg·°C) ไม่คิดการสูญเสียความร้อน \hfill \textbf{\small (Main)}
\begin{enumerate}[label=)]
\item ก) ประมาณ 26°C
\item ข) ประมาณ 40°C
\item ค) ประมาณ 60°C
\item ง) ประมาณ 80°C
\end{enumerate}
\item ที่อุณหภูมิ 0°C แก๊สอุดมคติ 1 โมลมีปริมาตรเท่าใดที่ความดัน 1 atm \hfill \textbf{\small (Main)}
\begin{enumerate}[label=)]
\item ก) 2.24 L
\item ข) 22.4 L
\item ค) 224 L
\item ง) 0.224 L
\end{enumerate}
\item ทำไมเงินจึงนำความร้อนได้ดีกว่าไม้ \hfill \textbf{\small (Main)}
\begin{enumerate}[label=)]
\item ก) เพราะเงินเป็นของแข็ง
\item ข) เพราะอนุภาคในเงินอยู่ชิดกันและสั่นสะเทือนถ่ายพลังงานให้กันเอง
\item ค) เพราะเงินมีมวลมากกว่า
\item ง) เพราะเงินมีอุณหภูมิสูงกว่า
\end{enumerate}
\item แก๊สในภาชนะปิดมีปริมาตรคงที่ ถ้าอุณหภูมิเพิ่มจาก 27°C เป็น 127°C ความดันจะเพิ่มขึ้นกี่เท่า \hfill \textbf{\small (Main)}
\begin{enumerate}[label=)]
\item ก) 2 เท่า
\item ข) 4/3 เท่า
\item ค) 1.33 เท่า
\item ง) 127/27 เท่า
\end{enumerate}
\item การเปลี่ยนหน่วยจากองศาเซลเซียสเป็นเคลวิน ทำอย่างไร \hfill \textbf{\small (Extra)}
\begin{enumerate}[label=)]
\item ก) °C + 273
\item ข) °C - 273
\item ค) (°C × 273) + 1
\item ง) °C × 273
\end{enumerate}
\item สารใดต่อไปนี้มีความจุความร้อนจำเพาะสูงที่สุด \hfill \textbf{\small (Extra)}
\begin{enumerate}[label=)]
\item ก) ทองแดง
\item ข) เหล็ก
\item ค) น้ำ
\item ง) อะลูมิเนียม
\end{enumerate}
\item น้ำมวล 3 kg ต้องการเพิ่มอุณหภูมิจาก 10°C เป็น 50°C ต้องใช้ความร้อนเท่าไร \hfill \textbf{\small (Extra)}
\begin{enumerate}[label=)]
\item ก) 504,000 J
\item ข) 168,000 J
\item ค) 126,000 J
\item ง) 252,000 J
\end{enumerate}
\item ความร้อนแฝงของการหลอมเหลว (L) ของน้ำแข็งมีค่า 334,000 J/kg หมายถึงอะไร \hfill \textbf{\small (Extra)}
\begin{enumerate}[label=)]
\item ก) ต้องใช้ 334,000 J เพื่อเพิ่มอุณหภูมิน้ำแข็ง 1 kg ขึ้น 1°C
\item ข) ต้องใช้ 334,000 J เพื่อละลายน้ำแข็ง 1 kg ที่จุดหลอมเหลว
\item ค) น้ำแข็ง 1 kg มีพลังงาน 334,000 J
\item ง) น้ำแข็งละลายได้ที่อุณหภูมิ 334,000°C
\end{enumerate}
\item การนำความร้อนผ่านแผ่นวัสดุหนา d พื้นที่ A ความนำความร้อน k ในเวลา t หาได้จากสูตรใด \hfill \textbf{\small (Extra)}
\begin{enumerate}[label=)]
\item ก) Q = kAdΔT/t
\item ข) Q = kAΔT t/d
\item ค) Q = kΔT/(Adt)
\item ง) Q = kdΔT/(At)
\end{enumerate}
\item แก๊สอุดมคติที่อุณหภูมิ 300 K มีความดัน 2 × 10⁵ Pa และปริมาตร 0.01 m³ มีกี่โมล \hfill \textbf{\small (Extra)}
\begin{enumerate}[label=)]
\item ก) 0.08 mol
\item ข) 0.8 mol
\item ค) 8 mol
\item ง) 80 mol
\end{enumerate}
\item ถ้าปริมาตรของแก๊สเพิ่มขึ้น 3 เท่าที่ความดันคงที่ อุณหภูมิจะเป็นอย่างไร \hfill \textbf{\small (Extra)}
\begin{enumerate}[label=)]
\item ก) เพิ่มขึ้น 3 เท่า
\item ข) ลดลง 3 เท่า
\item ค) คงที่
\item ง) เพิ่มขึ้น 9 เท่า
\end{enumerate}
\item ข้อใดไม่ใช่วิธีการถ่ายโอนความร้อน \hfill \textbf{\small (Extra)}
\begin{enumerate}[label=)]
\item ก) การนำ
\item ข) การพา
\item ค) การแผ่รังสี
\item ง) การเผาไหม้
\end{enumerate}
\item น้ำแข็ง 2 kg ที่ -10°C ละลายจนเป็นน้ำที่ 0°C ต้องใช้ความร้อนทั้งหมดเท่าไร (c น้ำแข็ง = 2,100 J/kg·°C, L = 334,000 J/kg) \hfill \textbf{\small (Extra)}
\begin{enumerate}[label=)]
\item ก) 668,000 J
\item ข) 709,200 J
\item ค) 42,000 J
\item ง) 726,200 J
\end{enumerate}
\item ภาชนะบรรจุแก๊ส 2 mol ที่ 2 atm และ 27°C มีปริมาตรเท่าไร (R = 0.082 L·atm/mol·K) \hfill \textbf{\small (Extra)}
\begin{enumerate}[label=)]
\item ก) 24.6 L
\item ข) 12.3 L
\item ค) 49.2 L
\item ง) 6.15 L
\end{enumerate}
\item จุดเดือดของน้ำที่ความดันบรรยากาศปกติ (1 atm) คือเท่าไร \hfill \textbf{\small (Extra)}
\begin{enumerate}[label=)]
\item ก) 0°C
\item ข) 50°C
\item ค) 100°C
\item ง) 212°C
\end{enumerate}
\item แก๊สชนิดหนึ่งมีปริมาตร 5 L ที่ความดัน 3 atm ถ้าความดันลดลงเหลือ 1 atm ที่อุณหภูมิเดิม ปริมาตรจะเป็นเท่าไร \hfill \textbf{\small (Extra)}
\begin{enumerate}[label=)]
\item ก) 15 L
\item ข) 5/3 L
\item ค) 5 L
\item ง) 8 L
\end{enumerate}
\item ทำไมพัดลมจึงช่วยทำให้อากาศเย็นลงได้ \hfill \textbf{\small (Extra)}
\begin{enumerate}[label=)]
\item ก) พัดลมสร้างความเย็นขึ้นมาเอง
\item ข) พัดลมพาอากาศร้อนออกไปและนำอากาศเย็นมาแทน
\item ค) พัดลมลดอุณหภูมิของอากาศโดยตรง
\item ง) พัดลมดูดความร้อนออกจากอากาศ
\end{enumerate}
\item เหล็ก 0.5 kg ที่ 80°C วางในน้ำ 1 kg ที่ 20°C อุณหภูมิสุดท้ายเป็นเท่าไร (c เหล็ก = 450, c น้ำ = 4,200 J/kg·°C) \hfill \textbf{\small (Extra)}
\begin{enumerate}[label=)]
\item ก) ประมาณ 25°C
\item ข) ประมาณ 35°C
\item ค) ประมาณ 50°C
\item ง) ประมาณ 60°C
\end{enumerate}
\item แก๊สอุดมคติ 0.5 mol ที่ 0°C และ 1 atm มีปริมาตรเท่าใด \hfill \textbf{\small (Extra)}
\begin{enumerate}[label=)]
\item ก) 5.6 L
\item ข) 11.2 L
\item ค) 22.4 L
\item ง) 44.8 L
\end{enumerate}
\item ข้อใดเกี่ยวกับการแผ่รังสีได้ถูกต้อง \hfill \textbf{\small (Extra)}
\begin{enumerate}[label=)]
\item ก) เกิดขึ้นเฉพาะในสุญญากาศ
\item ข) ต้องอาศัยตัวกลาง
\item ค) เป็นคลื่นแม่เหล็กไฟฟ้าที่เดินทางได้ทุกทิศทาง
\item ง) เกิดขึ้นเฉพาะจากดวงอาทิตย์
\end{enumerate}
\item กฎรวมแก๊ส (Combined Gas Law) คือข้อใด \hfill \textbf{\small (Extra)}
\begin{enumerate}[label=)]
\item ก) P₁V₁/T₁ = P₂V₂/T₂
\item ข) PV = nRT
\item ค) P ∝ V
\item ง) V ∝ T
\end{enumerate}
\item น้ำ 500 g ที่ 25°C ได้รับความร้อน 63,000 J อุณหภูมิสุดท้ายเป็นเท่าไร \hfill \textbf{\small (Extra)}
\begin{enumerate}[label=)]
\item ก) 30°C
\item ข) 40°C
\item ค) 55°C
\item ง) 70°C
\end{enumerate}
\item ถังแก๊สปริมาตร 20 L บรรจุแก๊สที่ 27°C ความดัน 5 atm ถ้าอุณหภูมิเพิ่มเป็น 127°C ที่ปริมาตรคงที่ ความดันจะเป็นเท่าไร \hfill \textbf{\small (Extra)}
\begin{enumerate}[label=)]
\item ก) 5.6 atm
\item ข) 6.67 atm
\item ค) 7.5 atm
\item ง) 10 atm
\end{enumerate}
\item ทำไมถ้วยกระดาษใส่น้ำร้อนไม่ลนมือ แต่ถ้าใส่เหล็กร้อนจะลน \hfill \textbf{\small (Extra)}
\begin{enumerate}[label=)]
\item ก) กระดาษมีความจุความร้อนจำเพาะต่ำ
\item ข) กระดาษนำความร้อนได้แย่ เลยไม่ส่งถ่ายความร้อนมาสู่มือเร็ว
\item ค) เหล็กมีความจุความร้อนจำเพาะสูง
\item ง) กระดาษดูดซับความร้อนไว้หมด
\end{enumerate}
\item หน่วยของความร้อนในระบบ SI คืออะไร \hfill \textbf{\small (Extra)}
\begin{enumerate}[label=)]
\item ก) แคลอรี (cal)
\item ข) จูล (J)
\item ค) องศาเซลเซียส (°C)
\item ง) บีทียู (BTU)
\end{enumerate}
\item แก๊สอุดมคติที่อุณหภูมิ 127°C ความดัน 2 atm มีปริมาตร 5 L มีกี่โมล \hfill \textbf{\small (Extra)}
\begin{enumerate}[label=)]
\item ก) 0.3 mol
\item ข) 0.4 mol
\item ค) 0.6 mol
\item ง) 1.2 mol
\end{enumerate}
\item ที่อุณหภูมิ -73°C แก๊สอุดมคติ 1 mol ที่ความดัน 1 atm มีปริมาตรเท่าใด \hfill \textbf{\small (Extra)}
\begin{enumerate}[label=)]
\item ก) 5.6 L
\item ข) 11.2 L
\item ค) 16.8 L
\item ง) 22.4 L
\end{enumerate}
\item กาแฟร้อนในถ้วยสามารถเย็นลงได้เร็วขึ้นโดยวิธีใด \hfill \textbf{\small (Extra)}
\begin{enumerate}[label=)]
\item ก) เป่าลมร้อนไปที่ถ้วย
\item ข) เป่าลมเย็นไปที่ถ้วย
\item ค) คลุมถ้วยด้วยฝาแน่น
\item ง) วางถ้วยในตู้เย็นปิด
\end{enumerate}
\item อะลูมิเนียม 1 kg ที่ 100°C วางในน้ำ 5 kg ที่ 20°C อุณหภูมิสุดท้ายเป็นเท่าไร (c อะลูมิเนียม = 900, c น้ำ = 4,200 J/kg·°C) \hfill \textbf{\small (Extra)}
\begin{enumerate}[label=)]
\item ก) ประมาณ 22°C
\item ข) ประมาณ 25°C
\item ค) ประมาณ 35°C
\item ง) ประมาณ 40°C
\end{enumerate}
\item กฎของอาโวกาโดรระบุว่าอะไร \hfill \textbf{\small (Extra)}
\begin{enumerate}[label=)]
\item ก) ที่ความดันและอุณหภูมิเดียวกัน แก๊ส 1 mol ทุกชนิดมีปริมาตรเท่ากัน
\item ข) P ∝ V ที่ T คงที่
\item ค) V ∝ T ที่ P คงที่
\item ง) P ∝ T ที่ V คงที่
\end{enumerate}
\item แก๊สอุดมคติในภาชนะ 5 L มีความดัน 3 atm อุณหภูมิ 300 K มีมวลโมเลกุลเฉลี่ย 40 g/mol มีกี่กรัม (R = 0.082 L·atm/mol·K) \hfill \textbf{\small (Extra)}
\begin{enumerate}[label=)]
\item ก) 24.6 g
\item ข) 49.2 g
\item ค) 73.8 g
\item ง) 12.3 g
\end{enumerate}
\item ข้อใดเป็นตัวนำความร้อนที่ดีที่สุด \hfill \textbf{\small (Extra)}
\begin{enumerate}[label=)]
\item ก) ไม้
\item ข) พลาสติก
\item ค) ทองแดง
\item ง) อากาศ
\end{enumerate}
\item น้ำแข็ง 1 kg ที่ 0°C ต้องใช้ความร้อนเท่าไรเพื่อเปลี่ยนเป็นไอที่ 100°C (L ของน้ำแข็ง = 334,000 J/kg, L ของการกลายเป็นไอ = 2,260,000 J/kg, c น้ำ = 4,200 J/kg·°C) \hfill \textbf{\small (Extra)}
\begin{enumerate}[label=)]
\item ก) 2,594,000 J
\item ข) 3,228,000 J
\item ค) 668,000 J
\item ง) 1,002,000 J
\end{enumerate}
\item แก๊สผสมประกอบด้วย Ne 2 mol และ He 1 mol ที่ความดันรวม 3 atm ความดันย่อยของ Ne เป็นเท่าไร \hfill \textbf{\small (Extra)}
\begin{enumerate}[label=)]
\item ก) 1 atm
\item ข) 1.5 atm
\item ค) 2 atm
\item ง) 3 atm
\end{enumerate}
\end{enumerate}
\end{document}
